\subsection{Problemática global de los residuos agrícolas y emisiones de GEI}

La gestión de residuos agrícolas constituye un desafío global con implicaciones directas en el cambio climático y la sostenibilidad ambiental. Los desechos orgánicos generados en este sector son una fuente significativa de emisiones de gases de efecto invernadero (GEI), como el metano (CH\textsubscript{4}),  producto de su descomposición anaeróbica \parencite{FAO2021}. El Quinto Informe de Evaluación del \textcite{IPCC2014} señala que el sector agrícola es responsable de aproximadamente el 24\% de las emisiones globales de GEI, siendo prácticas como la quema de residuos y el uso desmedido de fertilizantes factores agravantes de esta situación. Estas prácticas no solo incrementan las emisiones, sino que también contribuyen a la contaminación del suelo, el aire y el agua, afectando la biodiversidad y la salud humana. La magnitud de este problema ha impulsado la búsqueda de soluciones innovadoras y sostenibles que permitan una gestión eficiente de los residuos agrícolas y la reducción de sus impactos ambientales.

